\chapter{Sweat Test for Cystic Fibrosis}
\section*{What Is This Test?}
The sweat chloride test (pilocarpine iontophoresis sweat test) is the gold standard diagnostic test for cystic fibrosis (CF). CF is an autosomal recessive disorder caused by mutations in the CFTR gene (cystic fibrosis transmembrane conductance regulator, chromosome 7q31.2). This gene encodes a cAMP-regulated chloride channel. Defective CFTR function impairs chloride transport across epithelial cells. It produces thick, viscous secretions in the lungs (chronic bronchitis, bronchiectasis, recurrent respiratory infections with S. aureus and P. aeruginosa), pancreas (exocrine pancreatic insufficiency with fat malabsorption and steatorrhea, endocrine pancreatic insufficiency with CF-related diabetes), sweat glands (defective chloride reabsorption in the sweat duct produces high sweat chloride --- the basis of the sweat test), gastrointestinal tract (meconium ileus in neonates), male reproductive tract (congenital bilateral absence of the vas deferens [CBAVD], obstructive azoospermia), and hepatobiliary system (focal biliary cirrhosis). Pilocarpine iontophoresis stimulates sweat production. Sweat is collected on filter paper or in a Macroduct coil, and the chloride concentration is measured.
\section*{Alternative Names}
\begin{itemize}
  \item Sweat Chloride Test
  \item Pilocarpine Iontophoresis Sweat Test
  \item CF Sweat Test
  \item Sweat Electrolyte Test
\end{itemize}
\section*{Principle}
Pilocarpine iontophoresis: pilocarpine (a cholinergic agonist) is driven into the skin of the forearm by a small electrical current (iontophoresis) through positive and negative electrodes for 5 minutes. This stimulates maximal sweat production from the eccrine sweat glands. Sweat is collected for 30 minutes on pre-weighed filter paper (Gibson-Cooke method) or into a Macroduct coil. The collected sweat volume must be $\geq$75 mg (or $\geq$15 mcL for Macroduct) for a valid result. Sweat chloride concentration is measured by coulometric titration, chloride ion-selective electrode, or manual mercuric nitrate titration. Per CF Foundation guidelines, the test must be done at an experienced CF center (minimum 50--100 sweat tests annually) to ensure accuracy.
\section*{History}
In 1938, Dorothy Andersen first described cystic fibrosis as a distinct disease of childhood. She named it ``cystic fibrosis of the pancreas.'' In 1948, during a heat wave in New York, Paul di Sant'Agnese saw that children with CF developed salt-depletion heat prostration. In 1953, he and his colleagues reported abnormally high concentrations of sodium and chloride in the sweat of affected patients. In 1959, Lewis Gibson and Robert Cooke published the pilocarpine iontophoresis technique. It reliably stimulated sweat production and made the sweat chloride test the diagnostic standard. The Cystic Fibrosis Foundation issued diagnostic sweat testing guidelines in 1987 and updated them in 2007 and 2017. These set the current interpretive thresholds (<30, 30--59, and $\geq$60 mmol/L) and quality standards for testing sites. The CFTR gene was cloned in 1989. This enabled molecular confirmation and expanded newborn screening using immunoreactive trypsinogen followed by DNA analysis.
\section*{Why This Test Is Ordered}
\begin{itemize}
  \item Confirming the diagnosis of cystic fibrosis in patients with suggestive clinical features (recurrent pneumonia, chronic cough, steatorrhea, failure to thrive, or a positive newborn screening result)
  \item Diagnostic workup of children with recurrent or persistent respiratory infections, chronic sinusitis, nasal polyps, or bronchiectasis
  \item Evaluating pancreatic insufficiency, meconium ileus, or electrolyte abnormalities that suggest CF
  \item Investigation of unexplained male infertility due to congenital bilateral absence of the vas deferens
  \item Screening of siblings of affected children and other family members for genetic counseling
  \item Follow-up of a positive immunoreactive trypsinogen (IRT) newborn screening result to confirm or exclude CF
\end{itemize}
\section*{Primary vs Secondary Test}
The sweat chloride test is the primary, gold standard diagnostic test for cystic fibrosis. It is done after newborn screening (IRT followed by CFTR DNA analysis) or when clinical suspicion raises the possibility of CF. It gives definitive biochemical confirmation.
\section*{How This Test Is Performed}
The test is done on the volar surface of the forearm, or on the thigh in infants and older children. Pilocarpine, a cholinergic agonist, is applied to the skin. It is driven into the eccrine sweat glands by a low-intensity electrical current (iontophoresis) through electrodes for about 5 minutes. This causes no discomfort beyond a mild tingling sensation. Sweat is then collected for up to 30 minutes into pre-weighed filter paper (Gibson--Cooke method) or a Macroduct coil. A minimum volume (75 mg or 15 mcL) is required for a valid result. The chloride concentration is measured by coulometric titration or a chloride ion-selective electrode. Testing should be done at a Cystic Fibrosis Foundation-accredited center with trained personnel.
\section*{What Preparation Is Needed}
No fasting or special preparation is required, and no dietary changes are needed. The skin at the test site should be clean and free of lotions, creams, or ointments. These can interfere with sweat collection. Infants should be at least 2 weeks old and weigh at least 3 kg (about 6--7 pounds). Neonates may not produce enough sweat. Testing is typically deferred until 4--6 weeks of age when possible. The child should not be acutely dehydrated or febrile. The procedure takes about 45--60 minutes including collection and analysis.
\section*{Contraindications and Precautions}
\begin{itemize}
  \item There are no absolute contraindications. The test should be deferred in neonates younger than 2 weeks or weighing less than 3 kg. It should also be deferred in acutely ill or dehydrated children who may not produce enough sweat. Caution is required in infants with generalized edema, which can alter sweat electrolyte content. The electrical current used for iontophoresis is safe. It should not be applied to skin with open lesions or rashes. Electrodes must be placed on intact skin. A repeat test is required when results are intermediate (30--59 mmol/L) or when the first test is inconclusive.
\end{itemize}
\section*{Risks}
The risk is minimal. The iontophoresis current causes only a mild warming or tingling sensation. The skin may look slightly red for a few hours afterward. Sweat collection carries no risk of bleeding or infection and leaves no scar. Rarely, a child may become fussy during the 30-minute collection period.
\section*{What Happens After the Test}
Sweat chloride results are available the same day or within a few days. They are read as negative (<30 mmol/L), intermediate (30--59 mmol/L), or positive ($\geq$60 mmol/L). A positive result on two separate occasions confirms cystic fibrosis. So does one positive test with two disease-causing CFTR mutations. The family is then referred to a CF care center for counseling, genetic testing, and comprehensive management. Intermediate results prompt repeat testing, extended CFTR gene analysis, and functional studies such as nasal potential difference.
\section*{Understanding Results}
\subsection*{Negative (<30 mmol/L, <1 month: <30 mmol/L)}
CF is extremely unlikely (sweat chloride <30 mmol/L rules out CF in >99.9\% of cases). Normal CFTR function.
\subsection*{Intermediate (30--59 mmol/L)}
Equivocal --- cystic fibrosis cannot be excluded. Further testing is required: repeat sweat test, CFTR genetic testing, nasal potential difference (NPD) measurement, and extended CFTR gene analysis. CFTR-related disorder may present with intermediate sweat chloride.
\subsection*{Positive ($\geq$60 mmol/L)}
Diagnostic of cystic fibrosis. Per the CF Foundation diagnostic criteria, two positive sweat chloride tests $\geq$60 mmol/L on two separate occasions (or one test with two disease-causing CFTR mutations identified) confirm the diagnosis of CF. Sweat chloride levels in classical CF are typically 80--120+ mmol/L.
\section*{Normal Results}
Sweat chloride: <30 mmol/L (negative, CF excluded). 30--59 mmol/L (intermediate, equivocal, CF not excluded). $\geq$60 mmol/L (positive, CF diagnosed). Minimum sweat weight for a valid test: $\geq$75 mg (Gibson-Cooke) or $\geq$15 mcL (Macroduct). In neonates, sweat chloride $\geq$60 mmol/L is positive.
\section*{Follow-up Tests}
CFTR genetic testing (full gene sequencing for disease-causing variants, including the common delta-F508 deletion), immunoreactive trypsinogen (IRT) newborn screening, fecal elastase-1 (pancreatic insufficiency), nasal potential difference measurement (research/respecialized centers), and chest CT/pulmonary function tests.
\section*{References}
\begin{enumerate}
  \item Farrell PM, White TB, Ren CL, et al. Diagnosis of cystic fibrosis: consensus guidelines from the Cystic Fibrosis Foundation. J Pediatr. 2017;181S:S4--S15.
  \item Gibson LE, Cooke RE. A test for concentration of electrolytes in sweat in cystic fibrosis of the pancreas utilizing pilocarpine iontophoresis. Pediatrics. 1959;23(3):545--549.
  \item LeGrys VA, Yankaskas JR, Quittell LM, Marshall BC, Mogayzel PJ Jr. Diagnostic sweat testing: the Cystic Fibrosis Foundation guidelines. J Pediatr. 2007;151(1):85--89.
  \item di Sant'Agnese PA, Darling RC, Perera GA, Shea E. Abnormal electrolyte composition of sweat in cystic fibrosis of the pancreas. Pediatrics. 1953;12(5):549--563.
  \item Rowe SM, Miller S, Sorscher EJ. Cystic fibrosis. N Engl J Med. 2005;352(19):1992--2001.
  \item O'Sullivan BP, Freedman SD. Cystic fibrosis. Lancet. 2009;373(9678):1891--1904.
\end{enumerate}
\clearpage
\section*{Regulatory Status and Authorization Date}
This chapter describes a test category, not a specific commercial product. FDA, CDSCO, EU IVDR, and MHRA approval or registration dates therefore cannot be assigned without the assay manufacturer, product name, instrument, and intended use. For laboratory-developed tests, an individual agency approval date may not exist. See the Preface for the interpretation of regulatory status.

\clearpage
